\documentclass[11pt]{article}

\usepackage[margin=1in]{geometry}
\usepackage{amsmath,amssymb}
\usepackage[T1]{fontenc}

\title{Classic Reinforcement Learning: Core Mathematics}
\author{}
\date{}

\begin{document}

\maketitle

\section*{Overview}

Classic reinforcement learning (RL) is about an agent interacting with an environment to maximize long-term reward.

At time step $t$:
\begin{itemize}
    \item the agent observes a state $s_t$
    \item takes an action $a_t$
    \item receives a reward $r_{t+1}$
    \item transitions to the next state $s_{t+1}$
\end{itemize}

The goal is to learn a policy that maximizes expected cumulative reward.

\section{Markov Decision Process}

Classic RL is usually modeled as a Markov Decision Process (MDP):

\[
(\mathcal{S}, \mathcal{A}, P, R, \gamma)
\]

where:
\begin{itemize}
    \item $\mathcal{S}$ is the set of states
    \item $\mathcal{A}$ is the set of actions
    \item $P(s' \mid s,a)$ is the transition probability
    \item $R(s,a)$ or $R(s,a,s')$ is the reward function
    \item $\gamma \in [0,1]$ is the discount factor
\end{itemize}

The Markov property means the future depends only on the current state and action.

\section{Return}

The return is the discounted sum of future rewards:

\[
G_t = \sum_{k=0}^{\infty} \gamma^k r_{t+k+1}
\]

where:
\begin{itemize}
    \item $G_t$ is the return starting at time $t$
    \item $r_{t+k+1}$ is the reward received $k+1$ steps later
    \item $\gamma$ is the discount factor
\end{itemize}

Why it matters:
\begin{itemize}
    \item if $\gamma = 0$, the agent cares only about immediate reward
    \item if $\gamma \approx 1$, the agent cares more about long-term reward
\end{itemize}

\section{Policy}

A policy tells the agent what action to take.

For a stochastic policy:

\[
\pi(a \mid s) = \Pr(a_t = a \mid s_t = s)
\]

This means $\pi(a \mid s)$ is the probability of taking action $a$ in state $s$.

For a deterministic policy:

\[
a = \pi(s)
\]

This means the policy maps each state directly to one action.

\section{Objective Function}

The goal of RL is to maximize expected return:

\[
J(\pi) = \mathbb{E}_{\pi}[G_t]
\]

A more expanded form is:

\[
J(\pi) = \mathbb{E}_{\pi}\left[\sum_{t=0}^{\infty} \gamma^t r_{t+1}\right]
\]

This is the quantity the agent tries to optimize.

\section{State-Value Function}

The state-value function measures how good a state is under policy $\pi$:

\[
V^{\pi}(s) = \mathbb{E}_{\pi}[G_t \mid s_t = s]
\]

This is the expected return when starting from state $s$ and following policy $\pi$.

\section{Action-Value Function}

The action-value function, or Q-function, measures how good an action is in a state:

\[
Q^{\pi}(s,a) = \mathbb{E}_{\pi}[G_t \mid s_t = s, a_t = a]
\]

This is the expected return when starting in state $s$, taking action $a$, and then following policy $\pi$.

\section{Bellman Expectation Equation}

For the state-value function:

\[
V^{\pi}(s) = \sum_{a} \pi(a \mid s) \sum_{s'} P(s' \mid s,a)
\left[ R(s,a,s') + \gamma V^{\pi}(s') \right]
\]

This means:
\begin{itemize}
    \item value of the current state
    \item equals expected immediate reward
    \item plus discounted value of the next state
\end{itemize}

For the action-value function:

\[
Q^{\pi}(s,a) = \sum_{s'} P(s' \mid s,a)
\left[ R(s,a,s') + \gamma \sum_{a'} \pi(a' \mid s') Q^{\pi}(s',a') \right]
\]

This is the same recursive idea at the state-action level.

\section{Optimal Value Functions}

The optimal state-value function is:

\[
V^*(s) = \max_{\pi} V^{\pi}(s)
\]

The optimal action-value function is:

\[
Q^*(s,a) = \max_{\pi} Q^{\pi}(s,a)
\]

These represent the best possible long-term reward.

\section{Bellman Optimality Equation}

For the optimal state-value function:

\[
V^*(s) = \max_{a} \sum_{s'} P(s' \mid s,a)
\left[ R(s,a,s') + \gamma V^*(s') \right]
\]

For the optimal action-value function:

\[
Q^*(s,a) = \sum_{s'} P(s' \mid s,a)
\left[ R(s,a,s') + \gamma \max_{a'} Q^*(s',a') \right]
\]

These equations are central to value-based RL methods.

\section{Advantage Function}

The advantage function tells whether an action is better than average in a given state:

\[
A^{\pi}(s,a) = Q^{\pi}(s,a) - V^{\pi}(s)
\]

Interpretation:
\begin{itemize}
    \item if $A^{\pi}(s,a) > 0$, the action is better than average
    \item if $A^{\pi}(s,a) < 0$, the action is worse than average
\end{itemize}

\section{Temporal Difference Error}

The temporal difference (TD) error is:

\[
\delta_t = r_{t+1} + \gamma V(s_{t+1}) - V(s_t)
\]

This measures the gap between:
\begin{itemize}
    \item the current estimate $V(s_t)$
    \item a one-step improved target
\end{itemize}

\section{Q-Learning Update Rule}

One of the most important RL update rules is Q-learning:

\[
Q(s_t,a_t) \leftarrow Q(s_t,a_t) +
\alpha \left[
r_{t+1} + \gamma \max_{a'} Q(s_{t+1},a') - Q(s_t,a_t)
\right]
\]

where:
\begin{itemize}
    \item $\alpha$ is the learning rate
    \item $r_{t+1}$ is the immediate reward
    \item $\gamma$ is the discount factor
    \item $\max_{a'} Q(s_{t+1},a')$ is the estimated best future value
\end{itemize}

Core idea:
\begin{itemize}
    \item update the current Q-value
    \item toward reward plus best future estimate
\end{itemize}

\section{Policy Gradient}

A basic policy gradient formula is:

\[
\nabla_{\theta} J(\pi_{\theta}) =
\mathbb{E}_{\pi_{\theta}}
\left[
\nabla_{\theta} \log \pi_{\theta}(a_t \mid s_t) \, G_t
\right]
\]

where:
\begin{itemize}
    \item $\theta$ are the policy parameters
    \item the gradient increases probability of actions that led to high return
\end{itemize}

\section{Exploration vs. Exploitation}

A common strategy is $\epsilon$-greedy:

\[
\pi(a \mid s) =
\begin{cases}
1 - \epsilon + \frac{\epsilon}{|\mathcal{A}|}, & \text{if } a = \arg\max_{a'} Q(s,a') \\
\frac{\epsilon}{|\mathcal{A}|}, & \text{otherwise}
\end{cases}
\]

This means:
\begin{itemize}
    \item most of the time, choose the best-known action
    \item sometimes, explore a random action
\end{itemize}

\section{Compact Summary}

Classic RL is about maximizing expected discounted reward through interaction with an environment.

The core quantity is the return:

\[
G_t = \sum_{k=0}^{\infty} \gamma^k r_{t+k+1}
\]

From this we define:

\[
V^{\pi}(s) = \mathbb{E}_{\pi}[G_t \mid s_t = s]
\]

\[
Q^{\pi}(s,a) = \mathbb{E}_{\pi}[G_t \mid s_t = s, a_t = a]
\]

Bellman equations make RL recursive: value now equals immediate reward plus discounted future value.

\end{document}
